\documentclass[aspectratio=169,10pt]{beamer}

\usetheme{SimplePlus}
\usepackage[T1]{fontenc}
\usepackage[utf8]{inputenc}
\usepackage{amsmath,amssymb}
\usepackage{graphicx}
\usepackage{booktabs}
\usepackage{hyperref}

\graphicspath{{../self-material/book/part-1/images/}}

\title{Macroeconomics PhD Intro\\Lecture 2: The Solow Model}
\subtitle{Dynamics, Growth, Convergence, and Propagation}
\author{Swapnil Singh}
\institute{Lietuvos Bankas \and Kaunas University of Technology}
\date{Spring 2026}

\newcommand{\apxbutton}[1]{\hfill\hyperlink{#1}{\beamerbutton{Appendix}}}
\newcommand{\backbutton}[1]{\hfill\hyperlink{#1}{\beamerbutton{Back}}}

\begin{document}

\begin{frame}
  \titlepage
\end{frame}

\begin{frame}{Learning Goals}
\begin{itemize}
  \item Build the Solow model from assumptions to dynamics.
  \item Derive steady states, balanced growth, and convergence speed.
  \item Connect theory to facts on growth, returns, and factor shares.
  \item Use the model for impulse responses and log-linear analysis.
  \item Understand how the Solow framework serves as the core of virtually all modern business-cycle models.
\end{itemize}
\end{frame}

%% ============================================================
%% MOTIVATION
%% ============================================================

\begin{frame}{Motivation: The Puzzle of the Stable Capital-Output Ratio}
One of the central long-run economic facts is the stability of the capital-output ratio $K_t/Y_t$ over time: capital is roughly \textbf{three times annual GDP} in most advanced economies.

\medskip
Earlier growth theories (Harrod-Domar) took this stability as a \emph{technological} assumption---they simply built it into the model. But capital is only one production input. Production can be organized in many ways: highly capital-intensive or highly labor-intensive. Given the many possibilities, it is \textbf{not obvious} why $K_t/Y_t$ should be nearly constant at the macro level.

\medskip
\begin{block}{Solow's Starting Point}
Reconcile the stable capital-output ratio with the \emph{substitutability} of capital and labor. This reconciliation gave rise to:
\begin{enumerate}
  \item A framework for analyzing macroeconomic dynamics.
  \item Systematic methods for measuring technological change (growth accounting).
\end{enumerate}
\end{block}
\end{frame}

\begin{frame}{What the Solow Model Will Deliver}
The aggregate production function has three variables that can affect GDP growth: technology $A_t$, capital $K_t$, and labor $L_t$. The Solow model focuses on the \textbf{endogenous accumulation of capital} $K_t$.

\medskip
We will see that $K_t$ reacts not only to the saving rate $s$ but also to $A_t$ and $L_t$. After solving the model (which delivers a stable $K/Y$ ratio), we focus on two main takeaways:

\begin{enumerate}
  \item The fundamental source of long-run growth in per capita income is the growth in $A_t$.
  \item If all parameter values are common across economies, different countries converge to the same income per capita---both in level and growth rate---in the long run.
\end{enumerate}

\begin{block}{Road Map}
Basic model $\to$ steady state and dynamics $\to$ growing economy $\to$ prices and shares $\to$ convergence speed $\to$ business cycles and IRFs $\to$ log-linearization.
\end{block}
\end{frame}

\begin{frame}{Stylized Facts to Keep in Mind}
\begin{itemize}
  \item \textbf{Fact 1:} The capital-output ratio $K_t/Y_t$ is roughly stable over long horizons (approximately 3 in the U.S.).
  \item \textbf{Fact 2:} GDP per capita grows persistently at a roughly constant rate in advanced economies.
  \item \textbf{Fact 3:} The return to physical capital has been nearly constant over time.
  \item \textbf{Fact 4:} Factor shares (capital share $\approx 1/3$, labor share $\approx 2/3$) are roughly stable in the long run.
  \item \textbf{Fact 5:} There is conditional convergence across similar economies, but not strong unconditional convergence across all countries.
\end{itemize}

\medskip
The Solow model will account for all five of these facts with a single, parsimonious framework.
\end{frame}

%% ============================================================
\section{Basic Solow Model}
%% ============================================================

\begin{frame}{Technology and Assumptions}
We start with the simplest version: no technological progress, no population growth.

\begin{block}{Aggregate Production Function}
\[
Y_t = F(K_t,L_t)
\]
\end{block}
We interpret $Y_t$ as GDP. The function $F(K,L)$ satisfies:
\begin{enumerate}
  \item $F(K,L)$ is strictly increasing in both $K$ and $L$.
  \item $F(K,L)$ is strictly quasiconcave (strictly convex isoquants).
  \item Constant returns to scale: $F(cK,cL) = cF(K,L)$ for any $c>0$.
  \item $F(0,L)=0$: no output without capital.
  \item Inada conditions: $\lim_{K\to 0} F_1(K,L) = \infty$ and $\lim_{K\to \infty} F_1(K,L) = 0$.
\end{enumerate}
We also assume $F$ is twice continuously differentiable, so that FOCs generate continuous functions.
\end{frame}

\begin{frame}{Why These Assumptions Matter}
\begin{itemize}
  \item \textbf{CRS} allows reduction to a per-worker (or per-efficiency-worker) system---we collapse an $n$-dimensional problem into a one-dimensional difference equation.
  \item \textbf{Diminishing marginal product of capital} (from strict quasiconcavity) generates mean reversion of $k_t$: rich economies grow slower, poor economies grow faster.
  \item \textbf{Inada conditions} guarantee an interior steady state. They are stronger than strictly necessary---it suffices that $\lim_{k\to 0} f'(k) > \delta/s$ and $\lim_{k\to\infty} f'(k) < \delta/s$---but they simplify the exposition.
  \item \textbf{Strict quasiconcavity} gives well-behaved substitution between factors: isoquants are smooth and convex.
\end{itemize}
\begin{block}{Economic Logic}
At low $k$, the marginal product of capital is high, so investment return is high and capital grows. At high $k$, the marginal product is low and depreciation pressure dominates, so capital shrinks. The economy gravitates toward an intermediate steady state.
\end{block}
\end{frame}

\begin{frame}{What Does ``Neoclassical'' Mean Here?}
\begin{block}{Neoclassical Production Side}
\begin{itemize}
  \item Smooth substitution between capital and labor (convex isoquants).
  \item Diminishing marginal products for each input.
  \item Constant returns to scale at the aggregate level.
  \item Competitive pricing: factors paid their marginal products.
\end{itemize}
\end{block}
\begin{itemize}
  \item \textbf{Dynamic implication:} Strong convergence force through diminishing returns---the economy always returns to its steady state.
  \item \textbf{Empirical implication:} Disciplined mapping from factor shares and growth facts to structural parameters ($\alpha$, $\delta$, etc.).
  \item \textbf{Contrast with Harrod-Domar:} Fixed proportions (Leontief) technology implied knife-edge instability. The neoclassical substitutability of inputs removes this instability entirely.
\end{itemize}
\end{frame}

\begin{frame}{Per-Capita Representation}
\hypertarget{m-a1}{}
In the basic model, population and hours per capita are both constant and normalized to 1. Therefore:
\begin{itemize}
  \item Define per-worker (= per-capita) variables: $y_t \equiv Y_t/L_t$, $k_t \equiv K_t/L_t$.
  \item By CRS: $Y_t/L_t = F(K_t/L_t, 1)$.
  \item Define $f(k) \equiv F(k,1)$. Then:
\end{itemize}
\[
y_t = f(k_t) \tag{3.1}
\]
The intensive-form production function $f(\cdot)$ inherits all the properties of $F$: it is strictly increasing, strictly concave, and satisfies the Inada conditions in $k$.

\medskip
\textbf{Key point:} Because of the normalization $L_t=1$, per-capita and aggregate variables are the same in the basic model. The distinction will matter once we introduce population growth.
\apxbutton{a1}
\end{frame}

\begin{frame}{Capital Accumulation and Goods Market}
The second centerpiece of the Solow model is the equation describing the evolution of capital:
\begin{block}{Capital Accumulation}
\[
k_{t+1} = i_t + (1-\delta)k_t \tag{3.2}
\]
\end{block}
where $i_t$ is investment and $\delta \in (0,1)$ is capital's depreciation rate. The existing capital stock loses value $\delta k_t$ while being used.

\begin{block}{Goods Market Clearing}
\[
y_t = c_t + i_t \tag{3.3}
\]
\end{block}
Output is either consumed or invested. We implicitly assume goods are homogeneous---they can be used for both consumption and investment. Because output equals total income, total saving must equal total investment. (In an open economy with trade, this need not hold.)
\end{frame}

\begin{frame}{Saving Rule and Fundamental Equation}
\hypertarget{m-a2}{}
Instead of explicitly modeling consumption-saving decisions, the Solow model assumes consumers mechanically save a constant fraction of income:
\begin{block}{Behavioral Closure}
\[
i_t = s y_t, \quad s \in (0,1) \tag{3.4}
\]
\end{block}
This is the simplifying assumption that will be relaxed in Chapter 4, where saving arises from optimizing behavior.

\medskip
Inserting (3.4) into (3.2) and using (3.1):
\begin{block}{Fundamental Solow Equation}
\[
k_{t+1} = (1-\delta)k_t + s f(k_t) \tag{3.5}
\]
\end{block}
This difference equation is the \textbf{heart of the Solow model}. It expresses the entire dynamics of capital as a function of one endogenous variable $k_t$ and three exogenous objects: $\delta$, $s$, and $f$.
\apxbutton{a2}
\end{frame}

\begin{frame}{Economic Meaning of the Fundamental Equation}
\[
k_{t+1} = \underbrace{(1-\delta)k_t}_{\text{surviving capital}} + \underbrace{s f(k_t)}_{\text{new investment}}
\]

\begin{itemize}
  \item The only endogenous variable on the RHS is $k_t$. Given $k_0$, the entire future path $\{k_{t+1}\}_{t=0}^{\infty}$ is determined.
  \item Once we have the capital path, all other aggregates follow immediately:
\end{itemize}
\begin{block}{Immediate Implications}
\[
y_t=f(k_t),\qquad i_t=s f(k_t),\qquad c_t=(1-s)f(k_t)
\]
Transitional dynamics of output, investment, and consumption are all inherited from the path of $k_t$. This is the power of the one-dimensional state representation.
\end{block}
\end{frame}

%% ============================================================
\section{Steady State and Dynamics}
%% ============================================================

\begin{frame}{Steady State Definition}
A \textbf{steady state} is a situation where $k_t$ is constant over time: $k_t = \bar{k}$ for all $t$. From (3.5):
\begin{block}{Steady-State Condition}
\[
\bar{k} = (1-\delta)\bar{k} + s f(\bar{k})
\quad\Longrightarrow\quad
\delta \bar{k} = s f(\bar{k})
\]
\end{block}
\begin{itemize}
  \item In steady state, investment exactly replaces depreciation: $i = \delta \bar{k}$.
  \item All variables are constant: $\bar{y} = f(\bar{k})$, $\bar{c} = (1-s)f(\bar{k})$, $\bar{i} = sf(\bar{k})$.
  \item No growth in the basic model---the economy simply sits at $\bar{k}$ forever.
\end{itemize}
\end{frame}

\begin{frame}{Existence and Uniqueness of the Steady State}
\hypertarget{m-a3}{}
To see that a unique positive $\bar{k}$ exists, plot $\delta k$ and $sf(k)$ against $k$:
\begin{itemize}
  \item $\delta k$ is a straight line through the origin with slope $\delta$.
  \item $sf(k)$ is strictly increasing and strictly concave, reflecting diminishing returns.
  \item Near $k=0$: the Inada condition $f'(0)=\infty$ implies $sf(k)$ has infinite slope, so it starts above $\delta k$.
  \item As $k\to\infty$: the Inada condition $f'(\infty)=0$ implies $sf(k)$ eventually grows slower than $\delta k$.
  \item These two curves must cross exactly once at a positive $\bar{k}$.
\end{itemize}

\medskip
\textbf{Note:} The Inada conditions are stronger than necessary. It suffices that $\lim_{k\to 0} f'(k) > \delta/s$ (existence) and $\lim_{k\to\infty} f'(k) < \delta/s$ (uniqueness).
\apxbutton{a3}
\end{frame}

\begin{frame}{Cobb-Douglas Closed Form}
\hypertarget{m-a4}{}
A particularly useful production function is Cobb-Douglas: $F(K,L)=K^{\alpha}L^{1-\alpha}$, so $f(k)=k^{\alpha}$ with $\alpha\in(0,1)$. This satisfies all our assumptions, including Inada conditions.

\begin{block}{Analytical Steady State}
From $\delta\bar{k} = s\bar{k}^{\alpha}$:
\[
\bar{k}=\left(\frac{s}{\delta}\right)^{\frac{1}{1-\alpha}} \tag{3.6}
\]
\end{block}
\begin{itemize}
  \item $\bar{k}$ is increasing in $s$: higher saving $\Rightarrow$ more investment $\Rightarrow$ larger capital stock.
  \item $\bar{k}$ is decreasing in $\delta$: faster depreciation $\Rightarrow$ more capital lost each period $\Rightarrow$ lower steady state.
  \item Because $\bar{y}=f(\bar{k})=\bar{k}^{\alpha}$, steady-state output is also increasing in $s$ and decreasing in $\delta$.
\end{itemize}
\apxbutton{a4}
\end{frame}

\begin{frame}{Steady-State Capital-Output Ratio}
\hypertarget{m-a6}{}
Going back to our motivating fact: in the Cobb-Douglas case, the steady-state $k/y$ ratio is:
\begin{block}{Key Ratio}
\[
\frac{\bar{k}}{\bar{y}} = \frac{s}{\delta}
\]
\end{block}
\begin{itemize}
  \item In the basic model, $k_t/y_t$ is trivially constant in steady state (all variables are constant).
  \item But as we will see, even in a \emph{growing} economy where both $k_t$ and $y_t$ increase over time, this ratio eventually stabilizes---matching the data.
  \item $\bar{k}/\bar{y}$ is larger when $s$ is larger (more saving builds up capital relative to output).
  \item $\bar{k}/\bar{y}$ is smaller when $\delta$ is larger (faster depreciation erodes the stock).
  \item With $s\approx 0.2$ and $\delta\approx 0.05$, we get $\bar{k}/\bar{y}\approx 4$, broadly consistent with U.S.\ data.
\end{itemize}
\apxbutton{a6}
\end{frame}

\begin{frame}{Comparative Statics in Elasticity Form}
\hypertarget{m-a18}{}
\begin{block}{From $\bar{k}=(s/\delta)^{1/(1-\alpha)}$}
\[
\frac{\partial \ln \bar{k}}{\partial \ln s}=\frac{1}{1-\alpha},
\qquad
\frac{\partial \ln \bar{k}}{\partial \ln \delta}=-\frac{1}{1-\alpha}
\]
\end{block}
\begin{itemize}
  \item With $\alpha=1/3$: a 1\% increase in $s$ raises $\bar{k}$ by $1.5\%$.
  \item Output level response: $\partial \ln \bar{y}/\partial \ln s = \alpha/(1-\alpha)$.
  \item With $\alpha=1/3$: a 1\% increase in $s$ raises $\bar{y}$ by only $0.5\%$---diminishing returns dampen the output response.
  \item Higher curvature (lower $\alpha$) weakens sensitivity: with strong diminishing returns, the extra capital produces little extra output.
\end{itemize}
\apxbutton{a18}
\end{frame}

\begin{frame}{Global Dynamics: Cobweb Diagram}
\begin{center}
\includegraphics[scale=0.25]{17108034-9421-4dd7-a201-7af448ab5827-074_826_1066_631_552}
\end{center}
\end{frame}

\begin{frame}{Proof Sketch of Monotone Convergence}
\hypertarget{m-a5}{}
Suppose $k_t < \bar{k}$. We show:
\begin{enumerate}
  \item $k_{t+1} > k_t$: Because $k_t < \bar{k}$, we have $sf(k_t) > \delta k_t$ (the investment curve is above the depreciation line), so $k_{t+1} - k_t = sf(k_t) - \delta k_t > 0$.
  \item $k_{t+1} < \bar{k}$: Because $k_t < \bar{k}$ and the RHS of (3.5) is increasing in $k_t$, the value at $k_t$ is less than the value at $\bar{k}$, which equals $\bar{k}$.
\end{enumerate}

Therefore $k_t < k_{t+1} < \bar{k}$. Repeating: the sequence $\{k_t\}$ is monotone increasing and bounded above by $\bar{k}$.

\medskip
By the \textbf{Monotone Convergence Theorem}, the sequence has a limit. The limit must satisfy the steady-state equation, and since the steady state is unique, the limit is $\bar{k}$.

\medskip
The symmetric argument for $k_t > \bar{k}$ gives a monotone decreasing sequence bounded below by $\bar{k}$.
\apxbutton{a5}
\end{frame}

\begin{frame}{Convergence Intuition: Two Opposing Forces}
Rewrite the capital accumulation equation as:
\[
k_{t+1} - k_t = \underbrace{sf(k_t)}_{\text{investment force}} - \underbrace{\delta k_t}_{\text{depreciation force}}
\]
\begin{itemize}
  \item When $k_t$ is \textbf{small}: output per unit of capital $f(k_t)/k_t$ is large (diminishing returns). The constant saving rate implies large gross investment relative to the capital stock. Investment force wins.
  \item When $k_t$ is \textbf{large}: output per unit of capital is small. Gross investment cannot cover total depreciation $\delta k_t$. Depreciation force wins.
\end{itemize}

In growth-rate form:
\begin{block}{Growth Rate of Capital}
\[
\frac{k_{t+1}-k_t}{k_t} = \frac{s f(k_t)}{k_t} - \delta
\]
\end{block}
Since $f''(k)<0$ and $f(0)=0$, the average product $f(k)/k$ is decreasing in $k$. This generates the negative relationship between the growth rate and the level---the engine of convergence.
\end{frame}

\begin{frame}{Cobweb Method: General Setup}
\begin{block}{One-Dimensional Map}
\[
k_{t+1}=G(k_t), \quad G(k)\equiv (1-\delta)k+s f(k)
\]
\end{block}
\begin{itemize}
  \item Steady states are fixed points: $k^*=G(k^*)$, i.e., where the map crosses the 45-degree line.
  \item \textbf{Cobweb iteration:} Given $k_t$, move vertically to $G(\cdot)$ to read off $k_{t+1}$, then move horizontally to the 45-degree line to place $k_{t+1}$ on the horizontal axis. Repeat.
  \item The geometry around intersections determines local dynamics: the slope $G'(k^*)$ tells us everything about local behavior.
  \item This technique generalizes to any one-dimensional dynamical system---not just Solow.
\end{itemize}
\end{frame}

\begin{frame}{Local Cobweb Stability: Derivation}
\hypertarget{m-a17}{}
Define the deviation $x_t\equiv k_t-k^*$. First-order Taylor expansion around $k^*$:
\[
G(k_t)\approx G(k^*)+G'(k^*)(k_t-k^*)
\]
Since $G(k^*)=k^*$ (fixed point), we get:
\[
x_{t+1}\approx G'(k^*)\,x_t
\]
This is a simple linear recursion. Its solution is $x_t = [G'(k^*)]^t \, x_0$, so:
\begin{itemize}
  \item If $|G'(k^*)|<1$: deviations shrink $\Rightarrow$ \textbf{locally stable}.
  \item If $|G'(k^*)|>1$: deviations grow $\Rightarrow$ \textbf{locally unstable}.
  \item If $|G'(k^*)|=1$: borderline case, higher-order terms matter.
\end{itemize}
The \emph{sign} of $G'(k^*)$ determines whether convergence is monotone or oscillatory.
\apxbutton{a17}
\end{frame}

\begin{frame}{Cobweb Case 1: Monotone Convergence ($0<G'(k^*)<1$)}
\begin{block}{Condition}
\[
0<G'(k^*)<1
\]
\end{block}
\begin{itemize}
  \item Deviations keep the same sign and shrink each period: $x_t \to 0$ monotonically.
  \item Trajectory approaches $k^*$ from one side without overshooting.
  \item The \textbf{baseline neoclassical Solow model} is always in this case near steady state, because $G'(\bar{k}) = 1-\delta+sf'(\bar{k}) = 1 - \delta(1-\alpha) \in (0,1)$ for standard parameter values.
\end{itemize}
\end{frame}

\begin{frame}{Cobweb Case 2: Oscillatory Convergence ($-1<G'(k^*)<0$)}
\begin{block}{Condition}
\[
-1<G'(k^*)<0
\]
\end{block}
\begin{itemize}
  \item The sign of $x_t$ flips each period: the path overshoots the steady state repeatedly.
  \item But the absolute deviation shrinks, so oscillations are damped.
  \item Requires a locally \emph{decreasing} map around the fixed point---unusual for the standard Solow model, but can arise with non-standard production functions.
  \item Think of the bakery metaphor: too many ovens can reduce output if they crowd out workers and dough space.
\end{itemize}
\end{frame}

\begin{frame}{Cobweb Case 3: Instability ($|G'(k^*)|>1$)}
\begin{block}{Conditions}
\[
G'(k^*)>1 \quad \text{or} \quad G'(k^*)<-1
\]
\end{block}
\begin{itemize}
  \item Deviations expand over time: local divergence from the fixed point.
  \item If $G'(k^*)>1$: monotone explosion away from steady state.
  \item If $G'(k^*)<-1$: paths alternate in sign and explode---or settle into cycles or chaotic dynamics (forever non-monotone behavior that never settles to a repeated pattern).
  \item This is where multiple steady states and nonlinear attractors can emerge globally.
\end{itemize}
\end{frame}

%% ============================================================
\section{Other Dynamics}
%% ============================================================

\begin{frame}{Endogenous Growth Case}
\begin{columns}[T]
\column{0.45\textwidth}
\includegraphics[width=\linewidth]{17108034-9421-4dd7-a201-7af448ab5827-076_823_1043_436_509}
\column{0.55\textwidth}
Consider the case where $F_1(k,1)$ is uniformly above $\delta/s$. No positive steady state exists---$k_t$ grows without bound.
\begin{itemize}
  \item Special case: $y_t = Ak_t$ (linear in capital, $\alpha=1$ in Cobb-Douglas). Known as the ``$AK$ model.''
  \item Growth is ``endogenous'': it comes from within, not from exogenous technology.
  \item Two identical countries starting with different $k_0$ remain forever different---the gap in percentage terms stays constant.
  \item Since labor commands $\approx 2/3$ of income, the $AK$ setup is empirically implausible as a standalone model, but the concept is developed further in Chapter 13 with multiple accumulable inputs.
\end{itemize}
\end{columns}
\end{frame}

\begin{frame}{Poverty Traps}
\begin{columns}[T]
\column{0.4\textwidth}
\includegraphics[width=\linewidth]{17108034-9421-4dd7-a201-7af448ab5827-077_697_854_268_657}
\column{0.6\textwidth}
If $f(k)$ is not globally concave---e.g., there is a middle region with increasing returns (large infrastructure investments, building of transportation networks)---then the curve $k_{t+1}=G(k_t)$ can cross the 45-degree line multiple times.

\begin{itemize}
  \item In Figure 3.3: three steady states $\bar{k}_1 < \bar{k}_2 < \bar{k}_3$.
  \item $\bar{k}_1$ and $\bar{k}_3$ are \textbf{stable}; $\bar{k}_2$ is \textbf{unstable}.
  \item A country starting with $k_0 < \bar{k}_2$ converges to $\bar{k}_1$---the \textbf{poverty trap}.
  \item To escape: push $k_t$ above $\bar{k}_2$, e.g., by temporarily raising $s$ high enough that $\bar{k}_1$ and $\bar{k}_2$ disappear as intersections.
\end{itemize}
\end{columns}
\end{frame}

\begin{frame}{Complex Dynamics and Chaos}
\begin{columns}[T]
\column{0.4\textwidth}
\includegraphics[width=\linewidth]{17108034-9421-4dd7-a201-7af448ab5827-078_684_858_266_683}
\column{0.6\textwidth}
If $f(k)$ actually \emph{declines} at very high $k$ (more ovens than floor space in the bakery), then $G(k)$ can become decreasing. In this case:

\begin{itemize}
  \item Convergence, if it occurs at all, will not be monotone.
  \item $k_t$ can exhibit forever oscillating or even \textbf{chaotic} dynamics: great sensitivity to initial conditions, behavior that never settles to a repeated pattern.
  \item Mathematically: chaos requires $|G'(k^*)|>1$ with a non-monotone map.
  \item More esoteric in macroeconomics, but illustrates the full range of dynamics a one-dimensional system can produce.
\end{itemize}
\end{columns}
\end{frame}

\begin{frame}{Policy Interpretation of Non-Convex Cases}
\begin{itemize}
  \item \textbf{Multiple steady states} imply history dependence and threshold effects. Initial conditions determine long-run outcomes.
  \item A \textbf{temporary policy} that lifts $k_t$ above the unstable threshold can permanently alter the economy's trajectory---a ``big push'' argument for development policy.
  \item In practice, identifying true macro-level non-convexities (increasing returns at the aggregate level) has proved empirically difficult. The growth literature has not produced robust evidence for the kind of increasing returns needed.
\end{itemize}
\begin{block}{Caution}
These cases are useful for understanding possible mechanisms, but baseline neoclassical concavity remains the empirical benchmark.
\end{block}
\end{frame}

%% ============================================================
\section{Growing Economy}
%% ============================================================

\begin{frame}{Extending the Model: Why Growth?}
In the basic model, the long-run outcome is a steady state with \textbf{zero growth}. This is clearly counterfactual---GDP per capita grows persistently in the data.

\medskip
To address the growth facts from Chapter 2, we extend the model by allowing:
\begin{enumerate}
  \item The labor input $L_t$ to grow at rate $n$ (population growth and changes in hours worked).
  \item The technology level $A_t$ to grow at rate $\gamma$ (labor-augmenting technical change).
\end{enumerate}
\end{frame}

\begin{frame}{Harrod-Neutral Technical Change}
The extended production function takes the form:
\[
Y_t = F(K_t, A_t L_t)
\]
\begin{itemize}
  \item $A_t L_t$ is the total number of \textbf{efficiency units of labor} (or effective labor).
  \item Technology $A_t$ improves the labor input: this is \textbf{labor-augmenting} (Harrod-neutral) technical change.
  \item \textbf{Uzawa's Theorem (1961):} Labor-augmenting technical change is the \emph{only} form of technical progress consistent with exact balanced growth (where aggregate variables grow at constant rates). This is a deep theoretical result, not just a convenient assumption.
  \item Growth rate of $A_t$: $\gamma$. Growth rate of $L_t$: $n$.
  \item Note: $n$ has two origins---growing population and changes in hours worked per person (which has been declining in the long run).
\end{itemize}
\end{frame}

\begin{frame}{Transform to Efficiency Units}
\hypertarget{m-a7}{}
The same algebra as in the basic model yields $K_{t+1} = (1-\delta)K_t + sF(K_t, A_tL_t)$.

\begin{block}{Define Capital per Efficiency Unit}
\[
\tilde{k}_t \equiv \frac{K_t}{A_t L_t}
\]
\end{block}
Dividing both sides by $A_tL_t$ and noting that $K_{t+1}/(A_tL_t) = (1+\gamma)(1+n)\tilde{k}_{t+1}$:
\begin{block}{Fundamental Equation with Growth}
\[
(1+\gamma)(1+n)\tilde{k}_{t+1}=(1-\delta)\tilde{k}_t + s f(\tilde{k}_t) \tag{3.7}
\]
\end{block}
After normalization by $A_tL_t$, we recover a difference equation very similar to (3.5). The factor $(1+\gamma)(1+n)$ on the LHS reflects the extra ``depreciation'' from growing efficiency units: maintaining a given $\tilde{k}$ is harder when $A$ and $L$ grow.
\apxbutton{a7}
\end{frame}

\begin{frame}{Balanced Growth Path}
\hypertarget{m-a8}{}
The concept corresponding to the steady state is the \textbf{balanced growth path} (BGP). Along the BGP, $\tilde{k}_t$ is constant and all aggregate variables grow at constant rates.

\medskip
Setting $\tilde{k}_{t+1}=\tilde{k}_t=\overline{\tilde{k}}$ in (3.7):
\[
[(1+\gamma)(1+n) - (1-\delta)]\overline{\tilde{k}} = s f(\overline{\tilde{k}})
\]

With Cobb-Douglas $f(\tilde{k})=\tilde{k}^{\alpha}$:
\begin{block}{BGP Capital per Efficiency Unit}
\[
\overline{\tilde{k}}=\left(\frac{s}{(1+\gamma)(1+n)+\delta-1}\right)^{\frac{1}{1-\alpha}} \tag{3.8}
\]
\end{block}
Notice that $\gamma$ and $n$ enter like additional depreciation: faster growth in $A$ and $L$ means each unit of untransformed capital must ``make up for'' more effective depreciation.
\apxbutton{a8}
\end{frame}

\begin{frame}{Long-Run Growth Rate of Income per Capita}
\hypertarget{m-a9}{}
Income per capita is:
\[
y_t = \frac{Y_t}{L_t} = f(\tilde{k}_t)\cdot A_t
\]
Along the BGP, $f(\tilde{k}_t) = f(\overline{\tilde{k}})$ is constant. Therefore:
\[
\frac{y_{t+1}-y_t}{y_t}=\frac{A_{t+1}-A_t}{A_t}=\gamma
\]

\begin{block}{Fundamental Result}
The growth rate of per capita income in the long run equals $\gamma$---the rate of technological progress. \textbf{No other parameter affects it.} Encouraging saving (raising $s$) does \emph{not} affect the long-run growth rate.
\end{block}

This is perhaps the most surprising and important prediction of the Solow model: the engine of sustained growth is exogenous technology, not capital accumulation.
\apxbutton{a9}
\end{frame}

\begin{frame}{Transition Dynamics and Convergence}
In the short run, $\tilde{k}_t$ changes over time, affecting growth:
\[
\frac{y_{t+1}-y_t}{y_t} = \frac{f(\tilde{k}_{t+1})}{f(\tilde{k}_t)}(1+\gamma) - 1
\]
\begin{itemize}
  \item If $\tilde{k}_t < \overline{\tilde{k}}$: $\tilde{k}$ is growing, so $f(\tilde{k}_{t+1})/f(\tilde{k}_t) > 1$ and per-capita growth \textbf{exceeds} $\gamma$. The economy is ``catching up.''
  \item If $\tilde{k}_t > \overline{\tilde{k}}$: $\tilde{k}$ is falling, so per-capita growth is \textbf{below} $\gamma$. The economy is ``winding down'' to its BGP.
\end{itemize}
\begin{block}{Convergence Prediction}
Economies below their own BGP grow faster; economies above it grow slower. Poor countries (relative to their own BGP) grow faster during transition---this is the convergence prediction of the Solow model.
\end{block}
Note: this is \emph{conditional} convergence---it requires countries to share the same parameter values.
\end{frame}

\begin{frame}{Level Effect vs Growth Effect of Saving}
What happens when $s$ permanently increases?
\begin{itemize}
  \item $\overline{\tilde{k}}$ rises (from Eq.\ 3.8), so the \textbf{level} of $y_t$ on the BGP is permanently higher.
  \item During the transition to the new BGP, per-capita growth temporarily exceeds $\gamma$ as $\tilde{k}_t$ rises toward its new, higher value.
  \item But once the economy reaches the new BGP, per-capita growth returns to $\gamma$.
\end{itemize}
\begin{block}{Core Distinction}
The Solow model gives strong \textbf{level effects} of saving, but \textbf{no permanent growth effect}. Only changes in $\gamma$ can permanently alter long-run growth.
\end{block}
\textbf{Quantitative example:} With our calibration, doubling $s$ from 0.1 to 0.2 raises normalized output by about 17\%. The transition takes decades, temporarily boosting growth, but eventually the economy returns to growing at rate $\gamma$.
\end{frame}

%% ============================================================
\section{Prices and Shares}
%% ============================================================

\begin{frame}{Competitive Firm and Return to Capital}
\hypertarget{m-a10}{}
To connect the model to data on returns and factor shares, we need prices. Suppose firms maximize profit under competitive markets:
\begin{block}{Firm Problem}
\[
\max_{K_t,A_tL_t} F(K_t,A_tL_t)-r_tK_t-w_tA_tL_t \tag{3.9}
\]
\end{block}
Here, output is the num\'{e}raire (price = 1). Capital is owned by consumers and rented to firms at rate $r_t$; $w_t$ is the wage per efficiency unit of labor.

\medskip
From the FOC for $K_t$:
\[
r_t = F_1(K_t,A_tL_t) = f'(\tilde{k}_t)
\]
Along the BGP, $\tilde{k}_t$ is constant, so $r_t$ is constant---matching the data fact of a roughly stable return to capital.
\apxbutton{a10}
\end{frame}

\begin{frame}{Labor Share and Capital Share}
\hypertarget{m-a11}{}
The wage per efficiency unit satisfies: $w_t = F_2(K_t,A_tL_t)$.

\medskip
When $F$ has constant returns to scale, it is homogeneous of degree one. By \textbf{Euler's Theorem} for homogeneous functions:
\[
Y_t = K_t F_1(K_t,A_tL_t) + A_tL_t F_2(K_t,A_tL_t)
\]
Dividing by $Y_t$:
\[
1 = \underbrace{\frac{r_t K_t}{Y_t}}_{\text{capital share}} + \underbrace{\frac{w_t A_t L_t}{Y_t}}_{\text{labor share}}
\]
Capital share plus labor share equals one. Along the BGP, since both $r_t$ and $K_t/Y_t$ are constant, the capital share is constant. The labor share, being one minus the capital share, is also constant---matching the data.
\apxbutton{a11}
\end{frame}

\begin{frame}{Cobb-Douglas Benchmark for Factor Shares}
\begin{block}{If $Y_t=K_t^{\alpha}(A_tL_t)^{1-\alpha}$}
\begin{itemize}
  \item Capital share $= r_tK_t/Y_t = \alpha$.
  \item Labor share $= w_tA_tL_t/Y_t = 1-\alpha$.
  \item These are constant \textbf{even away from the BGP}---at any point on the transition path.
\end{itemize}
\end{block}
This is a special property of Cobb-Douglas. For more general CES production functions $Y=[\alpha K^{\rho}+(1-\alpha)(AL)^{\rho}]^{1/\rho}$, factor shares vary with $K/(AL)$ unless $\rho=0$ (Cobb-Douglas).

\medskip
\textbf{Key insight:} The Cobb-Douglas form is very convenient because it often simplifies algebra and leads to clean expressions. But this simplicity can be deceiving: it makes fundamental forces going in opposite directions cancel, rendering them invisible.
\end{frame}

\begin{frame}{Why Share Stability Matters Quantitatively}
\begin{itemize}
  \item Factor share data disciplines $\alpha$ directly: from U.S.\ national income accounts, $\alpha\approx 1/3$.
  \item With Cobb-Douglas, stable shares become a built-in structural feature, not an additional implication to check.
  \item The value of $\alpha$ pins down convergence speed $\lambda$ and all transition dynamics.
  \item The substitution elasticity between inputs is exactly 1 for Cobb-Douglas. Empirical estimates at the aggregate level rarely suggest major departures from 1, supporting the functional form choice.
\end{itemize}
\begin{block}{Empirical Caveat}
If measured shares trend strongly (as some recent work on declining labor shares suggests), Cobb-Douglas may require augmentation: human capital, markups, or a non-unit elasticity of substitution.
\end{block}
\end{frame}

%% ============================================================
\section{Convergence Speed and Evidence}
%% ============================================================

\begin{frame}{Local Linearization: No-Growth Case}
\hypertarget{m-a12}{}
To measure how fast the economy converges, linearize the fundamental equation around $\bar{k}$:
\[
\Delta k_{t+1}=\left(1-\delta+s f'(\bar{k})\right)\Delta k_t \tag{3.10}
\]
where $\Delta k_t \equiv k_t - \bar{k}$. Define the \textbf{convergence speed}:
\[
\lambda\equiv \delta-s f'(\bar{k}), \quad\text{so that}\quad
\frac{k_{t+1}-\bar{k}}{\bar{k}}=(1-\lambda)\frac{k_t-\bar{k}}{\bar{k}}
\]
A larger $\lambda$ means faster convergence (the gap shrinks more each period).

\medskip
For Cobb-Douglas, using the steady-state condition:
\[
\lambda=\delta(1-\alpha) \tag{3.11}
\]
Convergence is faster when $\alpha$ is small (stronger diminishing returns) and $\delta$ is large. Note: $s$ does \emph{not} affect $\lambda$ in the Cobb-Douglas case---two opposing forces cancel exactly.
\apxbutton{a12}
\end{frame}

\begin{frame}{Why $s$ Doesn't Affect $\lambda$ in Cobb-Douglas}
This deserves explanation because it is a subtle cancellation:
\begin{enumerate}
  \item For a given $k_t$, a larger $s$ means more investment, so $k_t$ has a \emph{larger} impact on $k_{t+1}$. This would \emph{slow} convergence (the map is steeper).
  \item But a larger $s$ also raises the steady-state $\bar{k}$, and at higher $\bar{k}$ the marginal product $f'(\bar{k})$ is smaller (diminishing returns). This makes the map \emph{flatter} near steady state, speeding up convergence.
\end{enumerate}

With Cobb-Douglas, these two forces exactly offset. For non-Cobb-Douglas production functions, $s$ generally does affect $\lambda$.

\medskip
\textbf{Lesson:} Cobb-Douglas simplicity can hide important economic forces. Always check whether a result is general or an artifact of the functional form.
\end{frame}

\begin{frame}{Local Linearization: Growth Case}
\hypertarget{m-a13}{}
With population and technology growth, linearizing (3.7) around $\overline{\tilde{k}}$:
\[
(1+\gamma)(1+n)\Delta\tilde{k}_{t+1}=\left(1-\delta+s f'(\overline{\tilde{k}})\right)\Delta\tilde{k}_t \tag{3.12}
\]
The convergence speed becomes:
\[
\lambda = 1-\frac{1-\delta+s f'(\overline{\tilde{k}})}{(1+\gamma)(1+n)}
\]
For Cobb-Douglas:
\[
\lambda=(1-\alpha)\left(1-\frac{1-\delta}{(1+\gamma)(1+n)}\right) \tag{3.13}
\]
Population growth and technology growth both increase $\lambda$ (speed up convergence), because they act like additional effective depreciation.
\apxbutton{a13}
\end{frame}

\begin{frame}{Speed Intuition}
\begin{center}
\includegraphics[width=0.72\linewidth]{17108034-9421-4dd7-a201-7af448ab5827-082_605_1528_1598_309}
\end{center}
\textbf{Figure 3.5:} A higher $\lambda$ means a flatter $G(\cdot)$ curve near the steady state, which geometrically means fewer cobweb steps to reach $\bar{k}$---faster convergence.
\end{frame}

\begin{frame}{From $\lambda$ to Half-Life}
The local convergence law is: gap$_{t+1} = (1-\lambda)\cdot$gap$_t$. After $h$ periods:
\[
\text{gap}_h = (1-\lambda)^h \cdot \text{gap}_0
\]
The \textbf{half-life} $h$ solves $(1-\lambda)^h = 1/2$:
\[
h = \frac{\ln(1/2)}{\ln(1-\lambda)} \approx \frac{0.693}{\lambda}
\]
\begin{itemize}
  \item With $\lambda = 0.049$ (model prediction): $h \approx 14$ years.
  \item With $\lambda = 0.02$ (empirical estimate): $h \approx 35$ years.
  \item The model predicts faster convergence than we observe---a key tension we will return to.
\end{itemize}
\end{frame}

\begin{frame}{Convergence in Data: All Countries vs.\ OECD}
\begin{columns}[T]
\column{0.5\textwidth}
\includegraphics[width=\linewidth]{17108034-9421-4dd7-a201-7af448ab5827-084_838_1073_275_537}

{\small Figure 3.6: All countries, 1960--2019. Source: PWT 10.0.}
\column{0.5\textwidth}
\includegraphics[width=\linewidth]{17108034-9421-4dd7-a201-7af448ab5827-085_871_1118_300_509}

{\small Figure 3.7: OECD countries, 1960--2019. Source: PWT 10.0.}
\end{columns}

\medskip
\begin{itemize}
  \item \textbf{All countries:} No systematic tendency for poor countries to grow faster---no unconditional convergence. But this does \emph{not} reject the Solow model, which predicts \emph{conditional} convergence.
  \item \textbf{OECD countries:} Clear negative relationship---poor countries in 1960 grew faster. Barro and Sala-i-Martin (1995) find similar patterns across U.S.\ states, Japanese prefectures, and European regions.
\end{itemize}
\end{frame}

\begin{frame}{Recent Convergence Patterns and Calibration}
\begin{columns}[T]
\column{0.5\textwidth}
\includegraphics[width=\linewidth]{17108034-9421-4dd7-a201-7af448ab5827-086_836_1078_275_519}

{\small Figure 3.8: All countries, 2000--2019. Source: PWT 10.0.}
\column{0.5\textwidth}
\begin{itemize}
  \item Kremer, Willis, and You (2022) show that unconditional convergence has emerged in recent decades.
  \item Their argument: policies, institutions, and human capital have become more similar across countries, making underlying parameters more homogeneous.
\end{itemize}
\end{columns}

\medskip
\textbf{Standard calibration:} $\alpha=1/3$, $\gamma=0.02$, $n=0.01$, $\delta\approx 0.046$. This gives $\lambda=0.049$, but the empirical estimate is $\lambda\approx 0.015$--$0.03$ (Barro and Sala-i-Martin, 2004). The model over-predicts convergence speed. Matching the data requires $\alpha\approx 0.73$, which can be rationalized by including human capital in the ``capital'' concept.
\end{frame}

%% ============================================================
\section{Business Cycles and IRFs}
%% ============================================================

\begin{frame}{Solow as Propagation Core}
The Solow model's usefulness goes far beyond growth theory. Due to its ability to account for the broad features of macroeconomic data, it constitutes the \textbf{core of virtually all modern business-cycle models}.

\begin{itemize}
  \item Business-cycle theories analyze how the economy reacts to shocks---the ``propagation mechanism.''
  \item Many cycle models are essentially Solow plus shocks: the state equation remains a modified version of (3.5).
  \item Three main types of shocks have been studied in the literature:
  \begin{enumerate}
    \item Neutral technology shocks (RBC tradition).
    \item Investment-specific technology shocks.
    \item Demand disturbances (Keynesian tradition).
  \end{enumerate}
\end{itemize}
\end{frame}

\begin{frame}{Shock 1: Neutral Technology (RBC Model)}
Production function: $y_t = A_t F(k_t, \ell_t)$ with variable labor $\ell_t$.
\begin{itemize}
  \item Shock: $A_t$ switches between high and low values. Economy moves toward different steady states.
  \item With constant $s$ and $\ell$, the model cannot match key business-cycle facts: (i) $\ell_t$ comoves with the cycle, (ii) $i_t$ is more volatile than $y_t$, (iii) $c_t$ is less volatile than $y_t$.
  \item So we allow $s$ and $\ell$ to react to $(k_t, A_t)$:
\end{itemize}
\[
k_{t+1}=(1-\delta)k_t+s(k_t,A_t)\,A_t\,F(k_t,\ell(k_t,A_t))
\]
\textbf{Note:} Under Cobb-Douglas, the distinction between Harrod-neutral ($F(K, AL)$) and Hicks-neutral ($AF(K,L)$) is inessential: $K^{\alpha}(AL)^{1-\alpha} = A^{1-\alpha}K^{\alpha}L^{1-\alpha}$.
\end{frame}

\begin{frame}{Shock 2: Investment-Specific Technology}
Goods market clearing becomes $y_t = c_t + i_t/\nu_t$, where $\nu_t$ captures the relative efficiency of producing investment goods:
\begin{itemize}
  \item When $\nu_t$ is high, it is cheaper to produce investment goods---one unit of consumption goods yields $\nu_t$ units of investment.
  \item This reflects a two-sector structure: the investment-goods sector has its own productivity.
\end{itemize}
\[
k_{t+1}=s\nu_tF(k_t,\ell)+(1-\delta)k_t
\]
\end{frame}

\begin{frame}{Shock 3: Demand-Driven Disturbance}
A Keynesian-flavored model: $c_t$ is exogenous and moves over time (the demand shock).
\begin{itemize}
  \item Total demand: $c_t + i_t = c_t/(1-s)$ (since $i_t/c_t = s/(1-s)$).
  \item When demand falls short of full-capacity output $F(k_t, \ell)$, total output is demand-determined and a fraction $u_t$ of the labor force becomes unemployed.
  \item $u_t = u(c_t, k_t)$ from: $F(k_t, \ell(1-u_t)) = c_t/(1-s)$.
\end{itemize}
\[
k_{t+1}=sF(k_t,\ell(1-u(c_t,k_t)))+(1-\delta)k_t
\]
This framework is Keynesian in spirit but begs the question of \emph{why} output can be below full capacity. Understanding the frictions behind demand shortfalls is central to later chapters.
\end{frame}

\begin{frame}{Impulse-Response Setup}
\hypertarget{m-a14}{}
In all three shock models, $k_{t+1}$ is a function of $k_t$ and a shock. To visualize the dynamics, we compute \textbf{impulse-response functions} (IRFs).

\medskip
\textbf{Setup:} Use the RBC version with Cobb-Douglas, exogenous $s$, and fixed $\ell$:
\[
k_{t+1}=sA_tk_t^{\alpha}\ell^{1-\alpha}+(1-\delta)k_t \tag{3.14}
\]
\begin{itemize}
  \item Before period 0: $A_t = \bar{A}$, $k_t$ is near $\bar{k}$.
  \item At $t=0$: $A_0 = (1+\varepsilon)\bar{A}$ (1\% shock).
  \item For $t\geq 1$: $A_t = (1+\rho^t\varepsilon)\bar{A}$, with $\rho=0.9$ (persistence).
  \item Starting from $k_0 = \bar{k}$, simulate (3.14) forward.
\end{itemize}
\apxbutton{a14}
\end{frame}

\begin{frame}{Impulse Response in the Solow-RBC Setup}
\begin{center}
\includegraphics[width=0.86\linewidth]{17108034-9421-4dd7-a201-7af448ab5827-091_581_1524_264_305}
\end{center}
\textbf{Figure 3.9:} Left: $A_t$ mean-reverts quickly. Right: $k_t$ responds gradually and persistently.

\begin{itemize}
  \item $k_t$ rises from $\bar{k}=9.066$ to a peak of 9.089 ($+0.25\%$), then gradually returns---a hump-shaped response.
  \item While $A_t$'s deviation fades by period $\sim$50, capital's response persists much longer.
  \item The peak capital deviation (0.25\%) is much smaller than the initial shock (1\%)---capital acts as a buffer.
  \item Eventually $k_t$ returns to $\bar{k}$: the convergence force is active even with shocks.
\end{itemize}
\end{frame}

%% ============================================================
\section{Log-Linearization}
%% ============================================================

\begin{frame}{Why Log-Linearization?}
\begin{itemize}
  \item Converts nonlinear dynamics into tractable \textbf{linear recursions} that can be solved analytically.
  \item Variables are expressed as \textbf{percent deviations} from steady state---natural economic units.
  \item Usually accurate for moderate shocks near the BGP. In our calibrated Solow model, the maximum approximation error is about $0.00001$ in units of $K_t$---indistinguishable from the exact solution when plotted.
  \item This technique is the workhorse of modern quantitative macroeconomics: DSGE models, RBC analysis, and New Keynesian models all rely heavily on log-linearization.
\end{itemize}
\end{frame}

\begin{frame}{Log-Deviation Definitions}
An arbitrary variable $x_t$ is expressed as:
\[
x_t=\bar{x}e^{\hat{x}_t},\qquad \hat{x}_t\equiv \log\left(\frac{x_t}{\bar{x}}\right) \tag{3.15}
\]
Since $\log(x_t/\bar{x}) \approx (x_t-\bar{x})/\bar{x}$ for small deviations (first-order Taylor expansion):
\[
\hat{x}_t\approx \frac{x_t-\bar{x}}{\bar{x}} \quad\text{(percent deviation from steady state)}
\]
And the exponential approximation:
\[
\bar{x}e^{\hat{x}_t}\approx \bar{x}(1+\hat{x}_t) \tag{3.16}
\]
These two approximations---replacing $e^{\hat{x}}$ with $1+\hat{x}$ and $\log(1+z)$ with $z$---are the building blocks of log-linearization.
\end{frame}

\begin{frame}{Linearized Capital Dynamics}
\hypertarget{m-a15}{}
Apply (3.15)--(3.16) to the Solow-RBC law (3.14):
\[
\bar{k}e^{\hat{k}_{t+1}}=s\bar{A}e^{\hat{A}_t}(\bar{k}e^{\hat{k}_t})^{\alpha}\ell^{1-\alpha}+(1-\delta)\bar{k}e^{\hat{k}_t}
\]
First-order approximation around zero deviations, using the steady-state identity $\delta\bar{k}=s\bar{A}\bar{k}^{\alpha}\ell^{1-\alpha}$:
\begin{block}{Log-Linearized Capital Law}
\[
\hat{k}_{t+1}=(1-\delta(1-\alpha))\hat{k}_t + \delta \hat{A}_t \tag{3.18}
\]
\end{block}
\begin{itemize}
  \item Let $\lambda\equiv \delta(1-\alpha)$---the same convergence speed as in (3.11).
  \item The log-linearization procedure is essentially the same as the linearization for convergence speed: both apply first-order Taylor approximation to (3.14).
  \item Capital responds to shocks with coefficient $\delta$ and has persistence $(1-\lambda)$.
\end{itemize}
\apxbutton{a15}
\end{frame}

\begin{frame}{Closed-Form Response and Output Deviation}
\hypertarget{m-a16}{}
With $\hat{A}_0=\varepsilon$ and $\hat{A}_t=\rho^t\varepsilon$ for $t\geq 1$, solving (3.18) forward:
\[
\hat{k}_{t+1}=\sum_{\tau=0}^{t}\rho^{t-\tau}(1-\lambda)^{\tau}\delta\varepsilon
= \rho^t \frac{1-\left(\frac{1-\lambda}{\rho}\right)^{t+1}}{1-\frac{1-\lambda}{\rho}}\,\delta\varepsilon
\]
Output deviation (from log-linearizing $y_t = A_t k_t^{\alpha}\ell^{1-\alpha}$):
\[
\hat{y}_t=\hat{A}_t+\alpha\hat{k}_t
\]
Output responds through two channels: the \textbf{direct effect} of the technology shock ($\hat{A}_t$) and the \textbf{propagated effect} through accumulated capital ($\alpha\hat{k}_t$). Log-linearization gives us these closed-form expressions---impossible to obtain from the nonlinear system.
\apxbutton{a16}
\end{frame}

\begin{frame}{Reading the IRF Formula}
\begin{itemize}
  \item The weights in the sum combine \textbf{shock persistence} $\rho$ and \textbf{capital persistence} $(1-\lambda)$.
  \item If $\rho>(1-\lambda)$: shock persistence dominates early dynamics. Capital builds up as long as the technology shock is active.
  \item If $\rho<(1-\lambda)$: capital's own persistence dominates. The hump in capital outlasts the shock itself.
  \item Larger $\lambda$ (faster convergence) shortens capital memory and compresses the hump-shaped response.
  \item Output has a direct term ($\hat{A}_t$, which decays as $\rho^t$) plus a propagated term ($\alpha\hat{k}_t$, which decays more slowly). This two-component structure is what generates the characteristic hump shape of output IRFs.
\end{itemize}
\end{frame}

\begin{frame}{Level-Linearization Alternative}
When the log formulation is inconvenient, we can linearize in levels:
\[
\Delta k_{t+1}=g_k\Delta k_t + g_A\Delta A_t
\]
where $g(k_t,A_t)\equiv(1-\delta)k_t+sA_tf(k_t)$, and $g_k$, $g_A$ are partial derivatives evaluated at steady state.

\medskip
For Cobb-Douglas:
\begin{itemize}
  \item $g_k = 1-\delta(1-\alpha)$ (same persistence as in log-linear version).
  \item $g_A = s^{1/(1-\alpha)}\delta^{-\alpha/(1-\alpha)}\bar{A}^{\alpha/(1-\alpha)}$.
\end{itemize}

The level-linearization and log-linearization give equivalent results for small deviations. The choice between them is one of convenience: log-deviations give percentage interpretations directly, while level-deviations may be easier when some variables can take negative values.
\end{frame}

%% ============================================================
\section{Wrap-up}
%% ============================================================

\begin{frame}{What the Solow Model Delivers}
\begin{enumerate}
  \item A disciplined dynamic core linking capital accumulation to growth, convergence, and the stability of key macroeconomic ratios.
  \item A clean mapping from parameters ($\alpha$, $s$, $\delta$, $\gamma$, $n$) to steady states, transition speeds, and long-run growth rates.
  \item Matches five stylized facts: stable $K/Y$, persistent per-capita growth, stable returns to capital, stable factor shares, and conditional convergence.
  \item A foundation for business-cycle propagation: virtually all modern macro models build on a version of the Solow framework.
  \item Analytical tractability via log-linearization, enabling closed-form impulse responses and quantitative analysis.
\end{enumerate}
\end{frame}

\begin{frame}{What the Solow Model Leaves Exogenous}
\begin{itemize}
  \item \textbf{Saving behavior:} The saving rate $s$ is assumed constant and exogenous. In reality, saving responds to interest rates, wealth, expectations about future income, etc.
  \item \textbf{Technology process:} The growth rate $\gamma$ of $A_t$ is taken as given. What drives technological progress? R\&D investment, institutions, human capital?
  \item \textbf{Labor supply:} In the basic model, labor is fixed. In reality, hours worked respond to wages, taxes, and the business cycle.
\end{itemize}

\begin{block}{Next Lecture (Chapter 4)}
Endogenous saving via \textbf{dynamic optimization}. Cass (1965) and Koopmans (1963) add microeconomic foundations: saving arises from rational, forward-looking consumers maximizing lifetime utility. This gives the same long-run implications as Solow, but allows for richer policy analysis and normative evaluation.
\end{block}
\end{frame}

%% ============================================================
\appendix
\section{Appendix Derivations}
%% ============================================================

\begin{frame}{A1. Aggregate to Per-Capita Form}
\hypertarget{a1}{}
\begin{enumerate}
  \item Start: $Y_t=F(K_t,L_t)$.
  \item Define $y_t\equiv Y_t/L_t$, $k_t\equiv K_t/L_t$.
  \item CRS implies $Y_t/L_t=F(K_t/L_t,1)$: divide both arguments by $L_t$.
  \item Define $f(k)\equiv F(k,1)$, so $y_t=f(k_t)$.
\end{enumerate}
Properties of $f$: Since $F$ is strictly increasing in $K$, so is $f$ in $k$. Since $F$ is CRS, $F_{11}<0$ (given strict quasiconcavity), so $f''<0$. The Inada conditions transfer: $f'(0)=\infty$, $f'(\infty)=0$.
\backbutton{m-a1}
\end{frame}

\begin{frame}{A2. Deriving the Fundamental Equation}
\hypertarget{a2}{}
\begin{enumerate}
  \item Capital law: $k_{t+1}=i_t+(1-\delta)k_t$.
  \item Saving rule: $i_t=s y_t$.
  \item Production: $y_t=f(k_t)$.
  \item Substitute (2) and (3) into (1): $k_{t+1}=(1-\delta)k_t+s f(k_t)$.
\end{enumerate}
This is a first-order nonlinear difference equation in $k_t$. Given $k_0>0$, it generates the entire time path of the economy.
\backbutton{m-a2}
\end{frame}

\begin{frame}{A3. Why Unique Positive Steady State?}
\hypertarget{a3}{}
Solve $\delta k = sf(k)$, equivalently $\delta = sf(k)/k = s\cdot\text{APK}(k)$.
\begin{itemize}
  \item LHS: constant $\delta > 0$.
  \item RHS: $s\cdot f(k)/k$. By Inada, $f'(0)=\infty \Rightarrow f(k)/k \to \infty$ as $k\to 0$.
  \item By concavity, $f(k)/k$ is strictly decreasing in $k$ (since $f''<0$ and $f(0)=0$).
  \item By Inada, $f'(\infty)=0 \Rightarrow f(k)/k \to 0$ as $k\to\infty$.
  \item A continuous strictly decreasing function from $\infty$ to $0$ crosses any positive horizontal line exactly once.
\end{itemize}
\backbutton{m-a3}
\end{frame}

\begin{frame}{A4. Derivation of Eq. (3.6)}
\hypertarget{a4}{}
With $f(k)=k^{\alpha}$:
\[
\delta\bar{k}=s\bar{k}^{\alpha}
\quad\Rightarrow\quad \bar{k}^{1-\alpha}=\frac{s}{\delta}
\quad\Rightarrow\quad \bar{k}=\left(\frac{s}{\delta}\right)^{\frac{1}{1-\alpha}}
\]
Steady-state output: $\bar{y} = \bar{k}^{\alpha} = (s/\delta)^{\alpha/(1-\alpha)}$.

Steady-state consumption: $\bar{c} = (1-s)\bar{y} = (1-s)(s/\delta)^{\alpha/(1-\alpha)}$.
\backbutton{m-a4}
\end{frame}

\begin{frame}{A5. Monotone Convergence Detail}
\hypertarget{a5}{}
\textbf{Case $k_t<\bar{k}$:}
\begin{itemize}
  \item $G(k_t) > k_t$ because $sf(k_t) > \delta k_t$ when $k<\bar{k}$. So $k_{t+1} > k_t$.
  \item $G(k_t) < G(\bar{k}) = \bar{k}$ because $G$ is increasing and $k_t < \bar{k}$. So $k_{t+1} < \bar{k}$.
  \item Combined: $k_t < k_{t+1} < \bar{k}$.
\end{itemize}

\textbf{Case $k_t>\bar{k}$:} Symmetric argument gives $\bar{k} < k_{t+1} < k_t$.

\medskip
In both cases, the sequence is monotone and bounded. By the Monotone Convergence Theorem, it has a limit $k^*$. Taking limits in $k_{t+1}=G(k_t)$ gives $k^*=G(k^*)$, and by uniqueness of the steady state, $k^*=\bar{k}$.
\backbutton{m-a5}
\end{frame}

\begin{frame}{A6. Derivation of $\bar{k}/\bar{y}=s/\delta$}
\hypertarget{a6}{}
From the steady-state condition $\delta\bar{k}=s f(\bar{k}) = s\bar{y}$:
\[
\frac{\bar{k}}{\bar{y}}=\frac{s}{\delta}
\]
This holds for \emph{any} production function satisfying our assumptions, not just Cobb-Douglas. The capital-output ratio in steady state depends only on the saving rate and the depreciation rate.
\backbutton{m-a6}
\end{frame}

\begin{frame}{A7. Derivation of Eq. (3.7)}
\hypertarget{a7}{}
Start from $K_{t+1} = (1-\delta)K_t + sF(K_t,A_tL_t)$. Divide both sides by $A_tL_t$:
\begin{align*}
\frac{K_{t+1}}{A_tL_t} &= (1-\delta)\frac{K_t}{A_tL_t}+s\frac{F(K_t,A_tL_t)}{A_tL_t}
\end{align*}
The LHS is:
\[
\frac{K_{t+1}}{A_tL_t} = \frac{A_{t+1}}{A_t}\cdot\frac{L_{t+1}}{L_t}\cdot\frac{K_{t+1}}{A_{t+1}L_{t+1}} = (1+\gamma)(1+n)\tilde{k}_{t+1}
\]
By CRS, $F(K_t,A_tL_t)/(A_tL_t) = F(K_t/(A_tL_t),1) = f(\tilde{k}_t)$. Therefore:
\[
(1+\gamma)(1+n)\tilde{k}_{t+1}=(1-\delta)\tilde{k}_t+s f(\tilde{k}_t)
\]
\backbutton{m-a7}
\end{frame}

\begin{frame}{A8. Derivation of Eq. (3.8)}
\hypertarget{a8}{}
Set $\tilde{k}_{t+1}=\tilde{k}_t=\overline{\tilde{k}}$ in (3.7):
\[
(1+\gamma)(1+n)\overline{\tilde{k}}=(1-\delta)\overline{\tilde{k}}+s\overline{\tilde{k}}^{\alpha}
\]
\[
[(1+\gamma)(1+n)-(1-\delta)]\overline{\tilde{k}}=s\overline{\tilde{k}}^{\alpha}
\]
\[
\overline{\tilde{k}}^{1-\alpha}=\frac{s}{(1+\gamma)(1+n)+\delta-1}
\]
\[
\overline{\tilde{k}}=\left(\frac{s}{(1+\gamma)(1+n)+\delta-1}\right)^{\frac{1}{1-\alpha}}
\]
\backbutton{m-a8}
\end{frame}

\begin{frame}{A9. Why Long-Run Per-Capita Growth Is $\gamma$}
\hypertarget{a9}{}
\[
y_t=\frac{Y_t}{L_t}=\frac{F(K_t,A_tL_t)}{L_t} = A_t\cdot\frac{F(K_t,A_tL_t)}{A_tL_t} = A_t\cdot f(\tilde{k}_t)
\]
On BGP, $\tilde{k}_t=\overline{\tilde{k}}$ is constant, so $f(\tilde{k}_t)$ is constant. Therefore:
\[
\frac{y_{t+1}-y_t}{y_t}=\frac{f(\overline{\tilde{k}})A_{t+1}-f(\overline{\tilde{k}})A_t}{f(\overline{\tilde{k}})A_t}=\frac{A_{t+1}-A_t}{A_t}=\gamma
\]
The $f(\overline{\tilde{k}})$ terms cancel, leaving only the growth rate of technology.
\backbutton{m-a9}
\end{frame}

\begin{frame}{A10. Deriving $r_t=f'(\tilde{k}_t)$}
\hypertarget{a10}{}
\begin{itemize}
  \item FOC of firm problem: $r_t=F_1(K_t,A_tL_t)$.
  \item From CRS: $f(\tilde{k}_t) \equiv F(\tilde{k}_t,1) = F(K_t,A_tL_t)/(A_tL_t)$.
  \item Differentiate both sides with respect to $K_t$:
\[
f'(\tilde{k}_t)\cdot\frac{1}{A_tL_t} = \frac{F_1(K_t,A_tL_t)}{A_tL_t}
\]
  \item Therefore $F_1(K_t,A_tL_t)=f'(\tilde{k}_t)$, i.e., $r_t = f'(\tilde{k}_t)$.
\end{itemize}
\backbutton{m-a10}
\end{frame}

\begin{frame}{A11. Factor Shares Under CRS}
\hypertarget{a11}{}
By Euler's theorem for homogeneous-of-degree-one functions:
\[
Y_t=K_tF_1(K_t,A_tL_t)+A_tL_tF_2(K_t,A_tL_t)
\]
Divide by $Y_t$:
\[
1=\frac{r_tK_t}{Y_t}+\frac{w_tA_tL_t}{Y_t}
\]
Hence labor share $= 1 -$ capital share. Both are constant along BGP since $r_t$ and $K_t/Y_t = \tilde{k}_t/f(\tilde{k}_t)$ are both constant when $\tilde{k}_t = \overline{\tilde{k}}$.

\medskip
Alternative: $w_t = f(\tilde{k}_t) - \tilde{k}_t f'(\tilde{k}_t)$, which is constant when $\tilde{k}_t$ is constant.
\backbutton{m-a11}
\end{frame}

\begin{frame}{A12. Deriving Eqs. (3.10)--(3.11)}
\hypertarget{a12}{}
Linearize $g(k)=(1-\delta)k+sf(k)$ at $\bar{k}$:
\[
\Delta k_{t+1}=g'(\bar{k})\Delta k_t=(1-\delta+s f'(\bar{k}))\Delta k_t
\]
Define $\lambda=\delta-s f'(\bar{k})$, so $g'(\bar{k}) = 1-\lambda$.

\medskip
For Cobb-Douglas: $f'(k)=\alpha k^{\alpha-1}$. At steady state, $s\bar{k}^{\alpha-1} = \delta/\alpha \cdot \alpha = \delta$ (from $s\bar{k}^{\alpha}=\delta\bar{k}$, so $sf'(\bar{k})=s\alpha\bar{k}^{\alpha-1}=\alpha\delta$). Therefore:
\[
\lambda=\delta-\alpha\delta=\delta(1-\alpha)
\]
\backbutton{m-a12}
\end{frame}

\begin{frame}{A13. Deriving Eqs. (3.12)--(3.13)}
\hypertarget{a13}{}
Linearize the growth-case map $(1+\gamma)(1+n)\tilde{k}_{t+1}=(1-\delta)\tilde{k}_t+sf(\tilde{k}_t)$ around $\overline{\tilde{k}}$:
\[
(1+\gamma)(1+n)\Delta\tilde{k}_{t+1}=(1-\delta+s f'(\overline{\tilde{k}}))\Delta\tilde{k}_t
\]
The convergence speed:
\[
\lambda=1-\frac{1-\delta+s f'(\overline{\tilde{k}})}{(1+\gamma)(1+n)}
\]
For Cobb-Douglas, $sf'(\overline{\tilde{k}}) = \alpha[(1+\gamma)(1+n)-1+\delta]$ (from the BGP condition). Substituting:
\[
\lambda=(1-\alpha)\left(1-\frac{1-\delta}{(1+\gamma)(1+n)}\right)
\]
\backbutton{m-a13}
\end{frame}

\begin{frame}{A14. IRF Calibration Inputs}
\hypertarget{a14}{}
In the chapter illustration:
\begin{itemize}
  \item $s=0.2$, $\delta=0.046$, $\alpha=1/3$.
  \item $\ell=1$, $\bar{A}=1$.
  \item Initial shock $\varepsilon=0.01$ (1\%), persistence $\rho=0.9$.
  \item Simulate Eq. (3.14) from $k_0=\bar{k}=(s\bar{A}/\delta)^{1/(1-\alpha)} = (0.2/0.046)^{1.5} \approx 9.066$.
  \item $A_t = (1+\rho^t\varepsilon)\bar{A}$ for $t=0,1,2,\ldots$
\end{itemize}
The implied convergence speed is $\lambda = 0.046 \times 2/3 \approx 0.031$.
\backbutton{m-a14}
\end{frame}

\begin{frame}{A15. Derivation of Eq. (3.18)}
\hypertarget{a15}{}
Start from the transformed law:
\[
\bar{k}e^{\hat{k}_{t+1}}=s\bar{A}e^{\hat{A}_t}(\bar{k}e^{\hat{k}_t})^{\alpha}\ell^{1-\alpha}+(1-\delta)\bar{k}e^{\hat{k}_t}
\]
Apply $e^{\hat{x}}\approx 1+\hat{x}$ and $(\bar{k}e^{\hat{k}_t})^{\alpha}\approx \bar{k}^{\alpha}(1+\alpha\hat{k}_t)$:
\[
\bar{k}(1+\hat{k}_{t+1})=s\bar{A}\bar{k}^{\alpha}\ell^{1-\alpha}(1+\hat{A}_t)(1+\alpha\hat{k}_t)+(1-\delta)\bar{k}(1+\hat{k}_t)
\]
Drop second-order terms ($\hat{A}_t\cdot\hat{k}_t\approx 0$) and use $\delta\bar{k}=s\bar{A}\bar{k}^{\alpha}\ell^{1-\alpha}$:
\[
\bar{k}\hat{k}_{t+1}=\delta\bar{k}(\hat{A}_t+\alpha\hat{k}_t)+(1-\delta)\bar{k}\hat{k}_t
\]
Divide by $\bar{k}$:
\[
\hat{k}_{t+1}=(1-\delta(1-\alpha))\hat{k}_t+\delta\hat{A}_t
\]
\backbutton{m-a15}
\end{frame}

\begin{frame}{A16. Closed-Form Responses}
\hypertarget{a16}{}
Given $\hat{A}_t=\rho^t\varepsilon$ and Eq. (3.18), iterate forward from $\hat{k}_0=0$:
\begin{align*}
\hat{k}_1 &= \delta\varepsilon \\
\hat{k}_2 &= (1-\lambda)\delta\varepsilon + \delta\rho\varepsilon \\
&\;\;\vdots \\
\hat{k}_{t+1} &= \sum_{\tau=0}^{t}\rho^{t-\tau}(1-\lambda)^{\tau}\delta\varepsilon
= \rho^t \frac{1-\left(\frac{1-\lambda}{\rho}\right)^{t+1}}{1-\frac{1-\lambda}{\rho}}\,\delta\varepsilon
\end{align*}
Output deviation: $\hat{y}_t=\hat{A}_t+\alpha\hat{k}_t = \rho^t\varepsilon + \alpha\hat{k}_t$.

With $\lambda=\delta(1-\alpha)$.
\backbutton{m-a16}
\end{frame}

\begin{frame}{A17. Cobweb Stability in Solow Form}
\hypertarget{a17}{}
For Solow, $G(k) = (1-\delta)k + sf(k)$, so:
\[
G'(k)=1-\delta+s f'(k)
\]
At $k^*$, local convergence requires $|G'(k^*)|<1$, i.e.:
\[
|1-\delta+s f'(k^*)|<1
\]
Since $1-\delta+sf'(k^*)>0$ for reasonable parameters, this simplifies to:
\[
1-\delta+s f'(k^*)<1 \quad\Leftrightarrow\quad sf'(k^*)<\delta \quad\Leftrightarrow\quad \lambda = \delta-sf'(k^*)>0
\]
In Cobb-Douglas: $\lambda=\delta(1-\alpha)$, so $0<\lambda<1$ for $\alpha\in(0,1)$ and any $\delta\in(0,1)$. Monotone convergence is guaranteed.
\backbutton{m-a17}
\end{frame}

\begin{frame}{A18. What Is an Elasticity?}
\hypertarget{a18}{}
\begin{block}{General Definition}
The \textbf{elasticity} of $y$ with respect to $x$ measures the percent change in $y$ per percent change in $x$:
\[
\varepsilon_{y,x} \;\equiv\; \frac{\partial y / y}{\partial x / x} \;=\; \frac{\partial y}{\partial x}\cdot\frac{x}{y} \;=\; \frac{\partial \ln y}{\partial \ln x}
\]
\end{block}
\textbf{Why use elasticities instead of plain derivatives?}
\begin{itemize}
  \item Elasticities are \textbf{unit-free}: they do not depend on whether capital is measured in dollars, euros, or physical units. A derivative $\partial \bar{k}/\partial s$ has units of ``capital per unit of saving rate,'' which is hard to interpret. An elasticity says: ``a 1\% rise in $s$ raises $\bar{k}$ by $\varepsilon\%$.''
  \item Elasticities give \textbf{proportional} comparisons: they tell you how large the response is \emph{relative to the starting point}, which is economically more meaningful than the absolute change.
  \item The $\partial\ln y / \partial\ln x$ representation is especially convenient because many economic relationships become linear in logs (e.g., Cobb-Douglas: $\ln y = \alpha \ln k$).
\end{itemize}
\end{frame}

\begin{frame}{A18 (cont.). Elasticity Applied to the Solow Steady State}
\hypertarget{a18b}{}
Start from $\bar{k}=(s/\delta)^{1/(1-\alpha)}$. Take logs:
\[
\ln \bar{k} = \frac{1}{1-\alpha}\left(\ln s - \ln \delta\right)
\]
Differentiate with respect to $\ln s$ (holding $\delta$ fixed):
\[
\frac{\partial \ln \bar{k}}{\partial \ln s} = \frac{1}{1-\alpha}
\]
\textbf{Interpretation:} A 1\% increase in the saving rate raises the steady-state capital stock by $1/(1-\alpha)$ percent.

\medskip
For output, $\ln \bar{y} = \alpha \ln \bar{k}$, so:
\[
\frac{\partial \ln \bar{y}}{\partial \ln s} = \alpha \cdot \frac{1}{1-\alpha} = \frac{\alpha}{1-\alpha}
\]
With $\alpha=1/3$: a 1\% rise in $s$ raises $\bar{k}$ by 1.5\% but $\bar{y}$ by only 0.5\%. Diminishing returns to capital attenuate the output response relative to the capital response.
\backbutton{m-a18}
\end{frame}

\end{document}
